\section{条件自举滤波器}\label{sec:bootstrap}

\subsection{滤波算法}\label{sec:alg}

假设已有一组独立近似服从时刻 $k - 1$ 后验概率密度函数 $p(\bm{x}_{k-1} \mid D_{k-1})$ 的随机样本 $\{\bm{x}_{k-1}(i) : i = 1, \dots, N\}$。条件自举滤波器是一种向前传播并更新这组样本的递推算法，旨在生成一组新样本 $\{\bm{x}_k(i) : i = 1, \dots, N\}$，使其近似服从时刻 $k$ 的后验密度 $p(\bm{x}_k \mid D_k)$。由此可见，该滤波器实质上是式~\eqref{eq:3} 与式~\eqref{eq:5} 递推关系的一种近似机械化仿真实现。

\emph{预测阶段}：将每一个后验样本代入系统状态方程，生成时刻 $k$ 的先验预测样本：
\begin{equation*}
\bm{x}_k^*(i) = \bm{f}_{k-1}\bigl(\bm{x}_{k-1}(i), \bm{w}_{k-1}(i)\bigr)
\end{equation*}
其中 $\bm{w}_{k-1}(i)$ 为从系统过程噪声概率密度函数 $p(\bm{w}_{k-1})$ 中独立抽取的样本。

\emph{更新阶段}：收到时刻 $k$ 的观测 $\bm{y}_k$ 后，评估每个先验样本对应的似然值，并为每个样本赋予归一化权重：
\begin{equation}\label{eq:8}
q_i = \frac{p\bigl(\bm{y}_k \mid \bm{x}_k^*(i)\bigr)}{\sum_{j=1}^N p\bigl(\bm{y}_k \mid \bm{x}_k^*(j)\bigr)}
\end{equation}
由此便在集合 $\{\bm{x}_k^*(i) : i = 1, \dots, N\}$ 上定义了一个离散概率分布，其中第 $i$ 个元素对应的点概率质量为 $q_i$。接着，从该离散分布中有放回地重采样 $N$ 次，生成新的后验样本集 $\{\bm{x}_k(i) : i = 1, \dots, N\}$，使得对任意 $j \in \{1, \dots, N\}$ 均满足
\begin{equation*}
\mathbb{P}\bigl(\bm{x}_k(j) = \bm{x}_k^*(i)\bigr) = q_i
\end{equation*}

上述预测与更新两个步骤共同构成递推滤波算法的一次完整迭代。为初始化递推过程，只需从已知初始先验密度 $p(\bm{x}_1)$ 中独立抽取 $N$ 个样本 $\bm{x}_1^*(i)$，并将其直接送入滤波器的更新阶段即可。经此递推产生的样本集 $\{\bm{x}_k(i)\}$ 便近似服从所需的后验条件密度 $p(\bm{x}_k \mid D_k)$。

\subsection{渐近合理性论证}\label{sec:just}

\emph{预测阶段}：算法预测步骤的合理性在直觉与数学上均十分直接。若 $\bm{x}_{k-1}(i)$ 为服从 $p(\bm{x}_{k-1} \mid D_{k-1})$ 的随机样本，且 $\bm{w}_{k-1}(i)$ 为与之独立并服从 $p(\bm{w}_{k-1})$ 的过程噪声样本，则由系统转移方程生成的 $\bm{x}_k^*(i) = \bm{f}_{k-1}(\bm{x}_{k-1}(i), \bm{w}_{k-1}(i))$ 恰好服从单步预测先验密度 $p(\bm{x}_k \mid D_{k-1})$。对全部 $N$ 个样本依次独立重复上述传播，即得一组独立同分布于 $p(\bm{x}_k \mid D_{k-1})$ 的先验样本 $\{\bm{x}_k^*(i) : i = 1, \dots, N\}$。

\emph{更新阶段}：更新步骤的理论依据源自 Smith 与 Gelfand \cite{r14} 证明的一项抽样重采样定理。他们指出，条件后验更新法则可以等价地实现为一种加权自举重采样。设已有一组独立取自连续密度函数 $G(\bm{x})$ 的样本 $\{\bm{x}_k^*(i) : i = 1, \dots, N\}$，而目标是从与 $L(\bm{x})G(\bm{x})$ 成正比的概率密度函数中抽取样本，其中 $L(\bm{x})$ 为已知非负函数。该定理表明：若从支撑集为 $\{\bm{x}_k^*(i) : i = 1, \dots, N\}$ 且各点概率质量为
\begin{equation*}
\frac{L\bigl(\bm{x}_k^*(i)\bigr)}{\sum_{j=1}^N L\bigl(\bm{x}_k^*(j)\bigr)}
\end{equation*}
的离散分布中抽取一个样本，则当样本规模 $N \to \infty$ 时，所抽样本的分布弱收敛于目标密度。只要将 $G(\bm{x})$ 对应为先验密度 $p(\bm{x}_k \mid D_{k-1})$，并将 $L(\bm{x})$ 对应为观测似然函数 $p(\bm{y}_k \mid \bm{x}_k)$，该定理便直接为式~\eqref{eq:8} 所刻画的更新过程提供了渐近合理性论证。

\subsection{讨论}\label{sec:disc}

如前所述，该采样滤波技术的核心优势在于，它既不对系统函数 $\bm{f}_k$ 与观测函数 $\bm{h}_k$ 的线性性或光滑性施加限制，也不对过程噪声与观测噪声的分布形态设限。算法运行的全部前置条件仅归结为以下三条：

\begin{itemize}
\item 初始先验密度 $p(\bm{x}_1)$ 可用于随机抽样；
\item 观测似然函数 $p(\bm{y}_k \mid \bm{x}_k)$ 具有可逐点计算的已知函数形式；
\item 系统过程噪声密度 $p(\bm{w}_k)$ 可用于随机抽样。
\end{itemize}

此外，基本算法结构极为简洁，极易编程实现。重采样更新阶段只需从区间 $(0, 1]$ 上的均匀分布抽取 $N$ 个独立随机数 $u_i$，对每个 $u_i$ 选取满足如下累计权重不等式的指标 $M$：
\begin{equation*}
\sum_{j=0}^{M-1} q_j < u_i \leq \sum_{j=0}^{M} q_j
\end{equation*}
并将对应的先验样本 $\bm{x}_k^*(M)$ 记作后验样本（其中约定 $q_0 = 0$）。对 $i = 1, \dots, N$ 重复此过程即完成整批重采样。由于各样本的预测传播与似然求值天然相互独立，该算法在大规模并行计算架构上的部署亦相当直接，从而为借助超大规模样本集实现实时非线性滤波创造了条件。

该算法的输出本身就是所需后验分布的一组样本，这在诸多应用中极为便利。例如，欲估计状态落入任意关注区域的后验概率，只需统计落入该区域的样本所占比例即可。同样，若已有充分理由确信后验分布为单峰形态，便可直接由样本分位数估计状态向量各分量的边际最高后验密度区域。状态后验均值、协方差乃至任意非线性状态函数的条件期望估计均可直接通过样本平均获得。

第~\ref{sec:just}~节 给出的条件自举滤波器合理性论证建立在 $N \to \infty$ 的渐近结论之上。对于有限样本规模，通常难以建立具有普适意义的严格误差界；同理，对于“究竟需要多少样本才能保证密度表示足以支撑滤波器稳定运行”这一关键工程问题，亦难以给出精确且可证的一般公式。不过可以确定的是，所需样本数 $N$ 至少受以下三个因素支配：

\begin{itemize}
\item 状态空间的维数；
\item 预测先验密度与观测似然函数之间典型的“重叠”程度；
\item 滤波器需要连续递推的时间步数。
\end{itemize}

先考察第一项因素：可以预期，样本规模 $N$ 会随状态空间维数迅速增长，其增长速率取决于状态向量各分量之间的统计耦合强度。在最温和的极限情形（即各状态分量彼此统计独立）下，所需的样本数并不会随空间维数增加。

$N$ 对后两项因素的依赖，则是重采样更新机制的直接后果。若似然函数 $p(\bm{y}_k \mid \bm{x}_k)$ 取显著值的状态空间区域相较于先验密度 $p(\bm{x}_k \mid D_{k-1})$ 的支撑区域狭小得多，那么绝大多数先验样本 $\bm{x}_k^*(i)$ 只会获得微乎其微的权重 $q_i$，从而在重采样中被淘汰。于是，远离似然高值区的先验样本实质上被浪费，而少数落入似然显著区附近的样本则会被反复抽中多次。当观测似然恰好落在先验密度较低的尾部区域时，这一退化现象会进一步加剧，因为该区域原本分布的先验样本就寥寥无几（另见文献 \cite{r14}）。经此反复筛选，后验概率密度函数的离散样本表示可能在短短数个时间步内迅速恶化；事实上，若系统不存在过程噪声，全部 $N$ 个样本甚至会迅速塌缩至单一状态点。原则上，可以凭借极大样本规模的“蛮力”策略缓解这一困境，但要面向具有实际意义的工程问题构建实用滤波器，仍需更为精巧的改进机制。为此，本文在第~\ref{sec:sim-bearings}~节 的仿真算例中提出并实现了对基本算法的两项修正：其一是\emph{样本抖动加糙}（roughening）过程，用以在重采样后对样本施加微弱独立扰动；其二是\emph{单步超前先验筛选}（prior editing）方案，用以增加落入显著似然附近的先验样本数量。在构造此类算法改进以进一步提升抽样效率方面，仍有广阔的研究空间。

Müller \cite{r12} 提出的抽样滤波算法在预测阶段与本文一致，但在更新阶段采用了不同策略：其更新过程是一种接受–拒绝抽样机制，依次考察每个先验样本 $\bm{x}_k^*(i)$，以正比于 $p(\bm{y}_k \mid \bm{x}_k^*(i))$ 的概率将其接受为后验样本，否则予以拒绝。该方案的一大缺点在于后验样本量是随机变量且会随时间递减。文献 \cite{r12} 的核心内容正是通过向原始先验样本拟合包络密度并从中重采样来扩充先验样本规模。
